\documentclass[11pt,a4paper]{article}
\usepackage[utf8]{inputenc}
\usepackage[margin=1in]{geometry}
\usepackage{hyperref}
\usepackage{array}
\usepackage{booktabs}
\usepackage{enumitem}
\usepackage{titlesec}
\usepackage{xcolor}

\definecolor{cudenvergold}{HTML}{8A724F}
\definecolor{darkgrey}{HTML}{333333}

\titleformat{\section}{\large\bfseries\color{darkgrey}}{\thesection}{1em}{}[{\color{cudenvergold}\titlerule[1.5pt]}]
\titleformat{\subsection}{\normalsize\bfseries\color{darkgrey}}{\thesubsection}{1em}{}

\hypersetup{
    colorlinks=true,
    linkcolor=darkgrey,
    filecolor=magenta,      
    urlcolor=cudenvergold,
    pdftitle={ISMG 2050 Fall 2026 Syllabus},
    pdfpagemode=FullScreen,
}

\begin{document}

\begin{center}
    {\LARGE \bfseries University of Colorado Denver} \\
    \vspace{2mm}
    {\Large \bfseries Business School} \\
    \vspace{4mm}
    {\color{cudenvergold}\hrule height 2pt}
    \vspace{3mm}
    {\large \bfseries ISMG 2050: Business Problem Solving Tools} \\
    {\large \bfseries Course Syllabus --- Fall 2026} \\
\end{center}

\vspace{0.5cm}

\section*{Course Logistics \& Instructor Information}
\begin{itemize}[leftmargin=*]
    \item \textbf{Instructor:} Franklin Diaz, Principal Engineer at Palo Alto Networks and Adjunct Faculty
    \item \textbf{Email:} \href{mailto:franklin.diaz@ucdenver.edu}{franklin.diaz@ucdenver.edu} (Canvas Inbox preferred)
    \item \textbf{Office Location:} Business School Building or Zoom (by appointment via Canvas Inbox)
    \item \textbf{Office Hours:} By appointment. Submit a request through the Canvas Inbox to schedule.
    \item \textbf{Student Success Coach (SSC):} Abdelkarim Mansour \\
          Coach Email: \href{mailto:abdelkarim.mansour@ucdenver.edu}{abdelkarim.mansour@ucdenver.edu}
    \item \textbf{Class Schedule:} Tuesday \& Thursday, 3:30 PM -- 4:45 PM
    \item \textbf{Classroom Location:} Business School Building, Room 2501
    \item \textbf{Term:} Fall 2026 (08/25/2026 -- 12/12/2026)
    \item \textbf{Delivery Mode:} In-person with asynchronous assignments via Canvas
\end{itemize}

\section*{Course Description}
ISMG 2050 focuses on the technology and problem-solving skills necessary for students to succeed both at school and in the business world. This course teaches how to make data-driven decisions using spreadsheets, data analysis, and visualization tools. Students solve problems in statistics, accounting, finance, marketing, management, and information systems. This hands-on course specifically utilizes Microsoft Excel and Tableau to provide students with practical problem-solving methods and tools necessary to succeed in the business community.

As a business core course, a grade of \textbf{C-} or better must be earned to satisfy Business graduation requirements and prerequisites for higher-level business courses.

\section*{Course Goals \& Learning Objectives}
By the end of this course, you will be able to:
\begin{itemize}
    \item Obtain basic proficiency in Microsoft Excel to provide a foundation for future business courses and professional development.
    \item Write formulas containing functions and operators in spreadsheet cells.
    \item Determine effective data display with formatting features, charts, and pivot tables.
    \item Write formulas for rules used in business decision-making situations.
    \item Understand introductory data management concepts and terminology.
    \item Obtain basic proficiency in Tableau to develop business intelligence capabilities.
    \item Transform raw data into meaningful visualizations and insights.
    \item Create interactive dashboards and stories, and handle multiple data sources in Tableau.
\end{itemize}

\section*{Required Texts \& Materials}
\begin{itemize}[leftmargin=*]
    \item \textbf{Primary Text:} Chamberlain, J. R. and Taylor, L. D., \textit{Spreadsheet Modeling for Business Decisions}, latest edition. Cengage Learning.
    \item \textbf{Software:} Microsoft Excel (installed locally), Tableau Public (free download at \url{https://public.tableau.com})
    \item \textbf{Additional Resources:} Course materials and datasets will be provided via Canvas modules each week.
\end{itemize}

\section*{Grading Criteria \& Weighting}
The finalized course grade is calculated based on the following weights:
\begin{center}
    \begin{tabular}{lc}
        \toprule
        \textbf{Item} & \textbf{Weight} \\
        \midrule
        Homework Assignments (Excel/Tableau) & 65\% \\
        Exams (Mid-Term and Final) & 20\% \\
        Tableau Final Project & 10\% \\
        Participation \& Engagement & 5\% \\
        \bottomrule
    \end{tabular}
\end{center}

\subsection*{Undergraduate Letter Grade Distribution}
\begin{center}
    \begin{tabular}{llp{7cm}}
        \toprule
        \textbf{Letter Grade} & \textbf{GPA Value} & \textbf{Percentage Range} \\
        \midrule
        A  & (4.0) & 93\% -- 100\% (Superior/Excellent) \\
        A- & (3.7) & 90\% -- 92.999\% \\
        B+ & (3.3) & 87\% -- 89.999\% \\
        B  & (3.0) & 84\% -- 86.999\% (Good/Better than Average) \\
        B- & (2.7) & 80\% -- 83.999\% \\
        C+ & (2.3) & 77\% -- 79.999\% \\
        C  & (2.0) & 74\% -- 76.999\% (Competent/Average) \\
        C- & (1.7) & 70\% -- 73.999\% (Minimum Passing for Business Core) \\
        D+ & (1.3) & 67\% -- 69.999\% \\
        D  & (1.0) & 64\% -- 66.999\% \\
        D- & (0.7) & 60\% -- 63.999\% \\
        F  & (0.0) & 0\% -- 59.999\% (Failing) \\
        \bottomrule
    \end{tabular}
\end{center}
\textit{Note: Grading policies of the CU Denver Business School state that the average GPA across all students in an undergraduate class should generally fall within the range of 2.3 (C+) to 3.0 (B). If necessary, grading scales may be slightly adjusted formulaically to align within this institutional guideline.}

\subsection*{Late Work Policy}
All assignments must be submitted through Canvas by the stated due date at \textbf{11:59 PM MST}. Assignments are subject to an immediate loss of \textbf{10 percentage points per day late}. No assignment will be accepted after 9 days late. On the 10th day, a score of \textbf{0} is recorded automatically.

\textbf{Mastery Retry:} If you receive a grade below C (74\%) on any homework worth 5\% or more of the course grade, you may retry that assignment for partial credit by meeting with me during office hours or by appointment. You must demonstrate mastery of the underlying concepts and submit within 7 days of receiving your original grade. The revised submission will be scored at up to 90\% of the original maximum.

\section*{Course Schedule \& Modular Breakdown}
The course is structured across 16 weeks, broken down into 8 core instructional content modules:

\begin{itemize}[leftmargin=*]
    \item \textbf{Module 1: Creating a Worksheet and a Chart (Weeks 1 \& 2)}
    \begin{itemize}
        \item \textit{Key Topics:} Microsoft Office ribbon components, Excel worksheet structure, simple functions/text/numbers, cell styling, and Pie Charts.
        \item \textit{Exercises:} Real Estate budget and chart project; Changing Values; Styles and Formatting; Loan Calculator; Comparing Laptops.
    \end{itemize}
    
    \item \textbf{Module 2: Formulas, Functions, and Formatting (Weeks 3 \& 4)}
    \begin{itemize}
        \item \textit{Key Topics:} Flash Fill, math formulas, workbook themes, Conditional Formatting, printing options.
        \item \textit{Exercises:} Klapore Engineering Salary Report; Cost Analysis; Kalto Security Outlet tracking; 4-Year College Cost Calculator; Cell Phone Service Summary.
    \end{itemize}

    \item \textbf{Module 3: Working with Large Worksheets, Charting, and What-If Analysis (Weeks 5 \& 6)}
    \begin{itemize}
        \item \textit{Key Topics:} Text rotations, advanced formatting, Absolute vs. Relative Cell References, IF functions, Sparkline charts, Format Painter, What-If Analysis, freezing panes.
        \item \textit{Exercises:} Manola Department Stores Six-Month Financial Projection; Logical Tests and Absolute Referencing; Fill Handle and Nested IF Statements; Census Data. \textbf{Quiz this week.}
    \end{itemize}

    \item \textbf{Module 4: Financial Functions, Data Tables, and Amortization (Weeks 7 \& 8)}
    \begin{itemize}
        \item \textit{Key Topics:} Financial Functions (PMT, PV, FV), Data Tables, Amortization Schedules, protecting/unprotecting worksheets.
        \item \textit{Exercises:} Cranfod Credit Union Mortgage Payment Calculator; Loan Payments; Retirement Planning; Down Payment Options.
    \end{itemize}

    \item \textbf{Module 5: Working with Multiple Worksheets and Workbooks (Weeks 9 \& 10)}
    \begin{itemize}
        \item \textit{Key Topics:} Formatting and consolidation techniques, Linear Series fills, custom format codes, custom cell styles, 3-D cell references.
        \item \textit{Exercises:} M\&S Provisions project; Consolidating Payroll; Custom Format Codes; Weather Data tracking.
    \end{itemize}

    \item \textbf{Module 6: Creating, Sorting, and Querying a Table (Weeks 11 \& 12) --- \underline{MID-TERM EXAM}}
    \begin{itemize}
        \item \textit{Key Topics:} Table creation/manipulation, calculated columns, XLOOKUP function, sorting, AutoFilter.
        \item \textit{Exercises:} Rating Bank Account Managers project; Table with Conditional Formatting; Advanced Functions; Converting Files. \textbf{Mid-Term Exam this week.}
    \end{itemize}

    \item \textbf{Module 7: Templates, Importing Data, Pivot Tables, and Macros (Weeks 13 \& 14)}
    \begin{itemize}
        \item \textit{Key Topics:} Excel Templates, SmartArt, image insertion, Pivot Tables, Pivot Charts, Macro basics.
        \item \textit{Exercises:} Mayor Insurance project; Template Consolidated Workbook; SmartArt Org Chart and Image; PivotTable and PivotChart creation.
    \end{itemize}

    \item \textbf{Module 8: Data Analysis, Charting, and Dashboards using Tableau (Weeks 15 \& 16) --- \underline{FINAL EXAM}}
    \begin{itemize}
        \item \textit{Key Topics:} Data management concepts, data visualization basics, Tableau data loading, building interactive dashboards.
        \item \textit{Exercises:} Bar Charts in Tableau; Data Analysis and Interpretation in Tableau; Course Review. \textbf{Tableau Final Project due. Final Exam this week.}
    \end{itemize}
\end{itemize}

\section*{AI Usage \& Academic Integrity Policy}

This course embraces the use of AI-assisted learning tools as part of your professional toolkit, with clear expectations about how and when they are appropriate.

\subsection*{Authorized Use (Permitted)}
\begin{itemize}[leftmargin=*]
    \item Decomposing nested formulas (IF/XLOOKUP) into manageable components
    \item Debugging M code logic in Power Query or VBA syntax errors
    \item Brainstorming visualization types for a given dataset
    \item Identifying the root cause of Excel error codes (\#VALUE!, \#REF!, etc.)
    \item Generating study questions for exam preparation
\end{itemize}

\subsection*{Unauthorized Use (Prohibited)}
\begin{itemize}[leftmargin=*]
    \item Submitting AI-generated completed worksheets or data files as your own work
    \item Having AI write full VBA macro scripts for direct submission
    \item Copying and pasting another student's work or sharing files with peers
    \item Using AI to complete in-class exams or quizzes (no AI permitted during exam periods)
\end{itemize}

\subsection*{AI Disclosure Requirement}
Every assignment submission must include an \textbf{AI Disclosure Statement} at the end of your document or embedded in your Canvas submission. Use this format:
\begin{verbatim}
AI Disclosure: I used [tool name, e.g., ChatGPT/Ollama/Gemini] for [brief description: e.g., "debugging VBA error codes"]. 
[OR] I did not use AI assistance on this assignment.
\end{verbatim}
Failure to disclose when AI was used constitutes academic dishonesty and will be treated as a submission of unauthorized work.

\subsection*{General Academic Honesty Statement}
Students are bound by the CU Denver Student Code of Conduct against plagiarism, cheating, fabrication, and falsification. Incorporating another's work into your submission without explicit authorization is academic dishonesty. Providing your Canvas login credentials to any other individual is considered academic dishonesty. Sharing files, copying formula text, or outsourcing graded dashboard setups will result in an automatic grade of ``F'' in the course and escalation to the Student Conduct Committee.

\section*{Strict Course Policies}

\subsection*{Best Mode of Communication}
All course-related correspondence must use the \textbf{messaging function within Canvas}. Canvas messages send an integrated copy to the instructor's inbox clearly marked with the course number. Do not email files directly from personal or student email programs; automated filters regularly block macro-embedded spreadsheets, and your communication may be lost.

Canvas Inbox is our secure, centralized, and auditable repository for all technical deliverables. From a data perspective, it ensures your work reaches its destination without being filtered by external systems.

\subsection*{Attendance \& Classroom Expectations}
Participation heavily influences mastery of technical tools. Attendance will be tracked on every session. If you are ill, notify the instructor via Canvas in advance.

Laptops are authorized strictly for active note-taking and working on in-class modules. Laptops shall \textbf{not} be used for social media, internet surfing, or messaging during class. \textbf{Mobile phones and texting are prohibited during class sessions.} Students who use their cell phones during active class will be asked to step out and will receive an unexcused absence for that day.

\subsection*{Visible ID Badge Requirement}
CU Denver policy requires all students, faculty, and staff in the Business School building to wear their official university ID badge visibly (lanyard, card holder, or clothing clip) at all times while inside the building.

\subsection*{Grades of Incomplete}
Incomplete (\textbf{I}) grades are granted strictly at the instructor's discretion for documented emergencies. An incomplete requires that the student has completed \textbf{at least 75\% of all coursework} with passing grades, and the remaining portion can be completed during the subsequent semester.

\section*{Student Success Coaching (SSC) Program}
This course features an integrated Student Success Coach, \textbf{Abdelkarim Mansour}, dedicated to supporting student performance.

The coach will proactively reach out if you miss a class, fail to submit a scheduled assignment, or if your Canvas average falls below a \textbf{``C'' grade}. 

If you receive a C or below on any major assignment (representing 5\% or more of the course grade), you are required to meet with the Success Coach before the next major assignment is due. During this session, the coach will assist you in understanding your mistakes and identifying specific steps for improvement. Once you complete this review meeting, you will receive \textbf{5\% of the points back} on that assignment (capped at 100\%).

\section*{University Resources \& Access Accommodation}
\begin{itemize}[leftmargin=*]
    \item \textbf{Disability Access (DRS):} Students entitled to academic adjustments must register with Disability Resources and Services in Student Commons Room 2116 (303-315-3510). Approved accommodation letters must be shared directly with the instructor.
    \item \textbf{Learning Resources Center (LRC):} Offers 24/7 online tutoring covering Excel/Tableau applications, plus in-person time management and study workshops (Student Commons \#1231, \url{https://www.ucdenver.edu/lrc}).
    \item \textbf{Mental Health \& Support:} The Student and Community Counseling Center (Tivoli Student Union \#454, 303-315-7270) provides cost-free, completely confidential psychological and mental health counseling.
    \item \textbf{Campus Safety (CARE Team):} Submit anonymous alerts about safety concerns at 303-315-7306 or via the online reporting form.
    \item \textbf{Title IX:} The university prohibits all forms of discrimination and sexual misconduct. Report concerns to the CARE Team at 303-315-7306.
\end{itemize}

\vfill
\noindent\rule{\textwidth}{0.4pt}
\small
\textit{Land Acknowledgement: We acknowledge the Cheyenne, Arapaho, and Ute nations as the original stewards of the land on which this university stands.}\\
\textit{This syllabus is a guide and may be adjusted as the term progresses. Any changes will be communicated promptly via Canvas. Please refer to the official CU Denver academic calendar for all term dates, reading days, and exam schedules.}

\end{document}
